\documentclass[12pt,a4paper]{article}

%% --- Encoding and fonts ---
\usepackage[utf8]{inputenc}
\usepackage[T1]{fontenc}
\usepackage{lmodern}

%% --- Language ---
\usepackage[english]{babel}

%% --- Mathematics ---
\usepackage{amsmath, amssymb, amsthm}
\usepackage{mathtools}

%% --- Layout ---
\usepackage[margin=1in]{geometry}
\usepackage{setspace}
\onehalfspacing

%% --- Figures and tables ---
\usepackage{graphicx}
\usepackage{booktabs}
\usepackage{caption}
\usepackage{subcaption}
\usepackage{float}

%% --- Appendix ---
\usepackage[title,titletoc]{appendix}

%% --- Cross-references and links ---
\usepackage{hyperref}
\hypersetup{colorlinks=true, linkcolor=blue, citecolor=blue, urlcolor=blue}

%% --- Bibliography ---
\usepackage{natbib}
\bibliographystyle{apalike}

%% --- Draft tools ---
\usepackage{xcolor}
\newcommand{\todo}[1]{\textcolor{red}{\textbf{[TODO: #1]}}}
\newcommand{\note}[1]{\textcolor{blue}{\textit{[Note: #1]}}}
% --- Custom Math Commands ---
\newcommand{\E}{\mathbb{E}}
\newcommand{\R}{\mathbb{R}}
\newcommand{\cA}{\mathcal{A}}          % asset space
\newcommand{\cH}{\mathcal{H}}          % productivity space
\newcommand{\cM}{\mathcal{M}}          % health space
\newcommand{\cF}{\mathcal{F}}          % formality space
\newcommand{\cQ}{\mathcal{Q}}          % pension-counter space
\newcommand{\ind}{\mathbf{1}}           % indicator function
\newcommand{\phij}{\phi_j}             % cohort weight
\newcommand{\jR}{j_R}                  % retirement age
% Type-specific shortcuts
\newcommand{\tL}{\theta_L}
\newcommand{\tH}{\theta_H}

%% --- Theorem environments ---
\newtheorem{assumption}{Assumption}
\newtheorem{proposition}{Proposition}
\newtheorem{definition}{Definition}
\newtheorem{remark}{Remark}

%% ====================================================================
\begin{document}

%% --- Title block ---
\title{\textbf{Spending Smarter under Demographic Pressure:}\\[6pt]
Fiscal Efficiency, Health, and Gender Dynamics\\
in Latin America\thanks{The work was financed with the support of the Latin America and the Caribbean Research Network of the Inter-American Development Bank. The opinions expressed in this publication are those of the authors and do not necessarily reflect the views of the Inter-American Development Bank, its Board of Directors, or the countries they represent.}}

\author{
Diego Ascarza-Mendoza\textsuperscript{1} \quad Hermilo Cort\'es\textsuperscript{1}\\[3pt]
Judith M\'endez\textsuperscript{1} \quad H\'ector J.\ Villarreal\textsuperscript{1}
\\[6pt]
\textsuperscript{1}Iniciativa para la Transici\'on Econ\'omica y Demogr\'afica (ITED),\\
Escuela de Gobierno y Transformaci\'on P\'ublica,\\
Tecnol\'ogico de Monterrey
}

\date{%
Prepared for the Inter-American Development Bank,\\
Fiscal Division\\[12pt]
\textit{Preliminary draft -- April 2026}\\
\textit{Please do not cite or circulate without permission}
}

\maketitle

%% --- Abstract ---
\begin{abstract}
Latin America and the Caribbean are undergoing an accelerated and heterogeneous demographic transition. Fertility has fallen below replacement in several countries, and longevity continues to rise, compressing the window of the demographic dividend and doubling old-age dependency ratios within two decades. This transition will impose unprecedented fiscal and social pressure on health, long-term care, and pension systems---precisely when fiscal space remains constrained, labor informality persists at high levels, and inequality remains entrenched.

This paper develops a stochastic overlapping-generations model with endogenous health investment, heterogeneous agents, and a pay-as-you-go pension system, calibrated for M\'exico, Costa Rica, and Panam\'a. Health capital is a continuous state variable that depreciates with age, affects survival probabilities, labor productivity, and the amenity value of consumption, and can be partially restored through medical spending subject to diminishing returns. The pension contribution rate adjusts endogenously to balance the system each period, creating a fiscal amplification loop: health investment that extends lives simultaneously improves productivity and raises the dependency ratio, with ambiguous net fiscal effects that only a calibrated general-equilibrium model can resolve. The model is used to study the fiscal consequences of demographic aging under policy inaction, evaluate a battery of tax and pension reform counterfactuals, and quantify the competing fiscal effects of labor formalization programs that expand the contributory base using observed income and health profiles from national household surveys. A calibration-based gender extension, exploiting sex-differentiated earnings trajectories and survival probabilities, decomposes fiscal sustainability pressures by sex within the existing equilibrium framework.
\end{abstract}

\noindent\textbf{Keywords:} demographic transition; fiscal sustainability; health expenditure; overlapping generations; labor informality; Latin America

\noindent\textbf{JEL codes:} E62, H51, I13, J11, O17

%% ====================================================================
%% SECTION 1 — INTRODUCTION
%% ====================================================================
\section{Introduction}
\label{sec:intro}

Latin America is aging. Yet it is doing so under conditions that distinguish the region sharply from the experience of high-income countries: fertility decline is rapid, formal labor markets remain incomplete, fiscal space is tight, and public health systems are chronically underfunded. The conjunction of these forces creates a structural dilemma that is at once demographic, fiscal, and institutional. The central question this paper addresses is how countries in the region can reallocate public health resources---spending \textit{smarter}---to meet the growing burden of aging populations without compromising fiscal sustainability or weakening the social security institutions on which workers depend.

The demographic challenge is well documented but its fiscal dimensions remain incompletely understood. By 2045, old-age dependency ratios in M\'exico, Costa Rica, and Panam\'a---our three country cases---are projected to at least double relative to their 2020 levels \citep{ECLAC2022}. Unlike the transitions observed in Western Europe or East Asia, this aging is occurring before these societies have achieved high-income status---a phenomenon often described as \textit{growing old before growing rich}. The implication is that the fiscal base from which health and pension commitments must be financed is narrower, and the informal sector, which contributes little to social security revenues, remains large. Critically, the aging that took Europe roughly 70 years to accomplish will take Latin America approximately 25 years, leaving far less time to legislate and implement the fiscal adjustments that demographic change demands \citep{WorldBankInformalizing2025}.

M\'exico illustrates the challenge in its starkest form. Between 2015 and 2024, our estimates show that Mexico's fiscal health gap widened from 3.0 to 4.8 percent of GDP: actual public health expenditure, held at roughly 2.8 percent of GDP, now covers barely more than a third of demographically projected requirements. Costa Rica and Panam\'a present a more favorable picture---both countries currently spend approximately in line with estimated needs---but demographic projections imply rising pressures in both cases as population aging accelerates over the next two decades. The cross-country variation in fiscal gaps, institutional architectures, and informality profiles is precisely why a comparative analysis of three countries is essential: it reveals which features of health-system design are structurally protective and which are artifacts of a particular country's fiscal position at a given moment. M\'exico's combination of high informality and a deteriorating fiscal gap identifies the most severe case; Costa Rica's historically generous contributory system through the CCSS (\textit{Caja Costarricense de Seguro Social}) shows the trajectory of a system approaching the limits of its fiscal surplus in health; and Panam\'a's institutional fragmentation between the CSS and the Ministry of Health illustrates efficiency losses from dual-system design. Together, the three cases span the relevant policy space and allow us to distinguish between country-specific and structural determinants of fiscal sustainability under aging.

A distinctive feature of the Latin American demographic transition, one that separates it analytically from high-income country experience, is the \textit{premature health crossing point}. In M\'exico, the proportion of the population in poor health crosses 50 percent at approximately age 60---a decade earlier than the 70--75 age range observed in high-income countries. This implies that the health system faces peak expenditure pressure precisely at the moment when workers are exiting the labor force and the contribution base is contracting. Skill differentials compound this pattern: among low-skilled workers, the health crossing point occurs even earlier, with rapid deterioration beginning in the late 40s, while high-skilled workers maintain good health status on average into their late 60s. This 20-year gap in healthy longevity between skill groups is the central empirical fact around which our model is built. Health is not one dimension of fiscal pressure among many---it is the first-order driver of the sustainability problem, and it is the analytical center of this paper.

We construct a stochastic overlapping-generations economy with endogenous health investment, heterogeneous agents who differ by skill type, and a government sector that operates a pay-as-you-go pension system. Health capital is a continuous state variable that depreciates with age and can be partially restored through medical spending, subject to diminishing returns. Health simultaneously affects survival probabilities, labor productivity, and the amenity value of consumption---making it a state variable that propagates through every margin of household decision-making. The model is solved as a stationary competitive equilibrium in which factor prices are determined endogenously through firm optimization and market clearing. We calibrate the model for M\'exico, Costa Rica, and Panam\'a using data from the Global Burden of Disease \citep{GBD2021}, the OECD Taxing Wages series \citep{OECDTaxingWages2026}, the WHO Global Health Expenditure Database \citep{WHOGHED2024}, and national household surveys.

The model's health formation technology generates a \textit{longevity channel} that operates as a macroeconomic transmission mechanism \citep{Carrasco2026}. When households invest in health, they live longer, which strengthens the life-cycle savings motive, raises the aggregate capital stock, and alters the equilibrium interest rate and wage. This channel operates asymmetrically across skill groups: low-skilled workers, who face earlier and steeper health depreciation, have shorter effective planning horizons and accumulate less precautionary savings, reinforcing the health-income divergence that the premature health crossing point generates. In high-income countries, where health deterioration begins later and proceeds more gradually, the longevity channel operates more uniformly across the skill distribution. In Latin America, it amplifies inequality.

The longevity channel interacts with the pension system through a fiscal amplification mechanism. The pension contribution rate $\tau^p$ adjusts endogenously to balance the system each period, equaling the replacement rate $\kappa$ scaled by the dependency ratio $N^R/N^W$---the ratio of retirees to workers. As the population ages and the dependency ratio rises, $\tau^p$ increases mechanically, widening the tax wedge on formal labor income. Health investment that extends productive working lives partially offsets this pressure by keeping workers in the labor force longer and by improving their productivity while active. But the offset is partial: healthier workers who survive into retirement also draw pension benefits for more periods, raising $N^R$. The net fiscal effect of improved population health is therefore ambiguous in sign and quantitatively important---precisely the kind of question that requires a calibrated general-equilibrium model to resolve.

The model's structure generates several classes of quantitative results. First, lifecycle profiles of health capital, consumption, and savings by skill type, where the interaction of age-dependent health depreciation with endogenous investment produces divergent trajectories: high-skilled households maintain health longer, save more through the longevity channel, and arrive at retirement with substantially larger asset buffers. Second, the fiscal gap decomposition---applying age-expenditure profiles to the model's stationary population distribution yields country-specific estimates of the divergence between demographically projected health needs and actual public spending. Third, the equilibrium pension contribution rate responds endogenously to the dependency ratio, feeding back into household savings and labor supply decisions. Fourth, unlike partial-equilibrium analysis, this model allows health-driven changes in aggregate savings to move the interest rate and the wage, so that a deterioration in population health affects all agents through factor prices, not only the households directly hit. Fifth, the government budget closure produces an implied debt-to-output ratio for each country's baseline demographic structure---a direct fiscal sustainability diagnostic that identifies where adjustment pressures will be most severe.

The baseline model is calibrated to the formal-sector population. To study the fiscal consequences of expanding the contributory base, we design a formalization experiment. The total population $A$ is partitioned into formal workers $B$ and informal workers $C$, with high-skilled workers overrepresented in $B$ and low-skilled workers overrepresented in $C$. The experiment simulates a formalization program that incorporates a fraction of $C$ into $B$, using observed income and health profiles from national household surveys---ENOE for M\'exico, ENAHO for Costa Rica, and the Encuesta de Hogares for Panam\'a. The skill composition of the newly formalized cohort---predominantly low-skilled, with earlier health depreciation and lower lifetime earnings---generates competing fiscal effects: a broader contribution base that reduces the equilibrium pension rate $\tau^p$, against higher per-capita health expenditure claims and lower average productivity. The model quantifies the net fiscal impact and identifies, for each country, the threshold formalization rate at which the contribution-base effect dominates the expenditure-cost effect.

Beyond the formalization experiment, we examine a battery of policy counterfactuals that vary the fiscal instruments and institutional design parameters of the health and pension systems. These include: (i)~expansion of public health coverage through contributory versus tax-financed channels, with explicit tracking of effects on the tax wedge and the equilibrium contribution rate; (ii)~preventive health interventions targeted at ages 25--35, where the model's lifecycle profiles suggest the welfare return is steepest; (iii)~modified pension vesting requirements that alter the cumulative contribution threshold; and (iv)~joint experiments that combine health and pension reforms to identify complementarities. The joint experiments are essential because these instruments interact: a reduction in the pension vesting threshold raises the value of formal employment for workers most likely to be deterred by high contribution rates, and this effect is amplified when combined with preventive health investments that keep those workers productive longer. Both independent and joint counterfactuals are reported, addressing directly the question of whether policy complementarities are quantitatively important.

The model's architecture supports a natural extension toward gender-differentiated analysis. Rather than constructing a separate household bargaining module, we exploit the fact that sex-differentiated income profiles and survival probabilities already embed the cumulative effects of labor market discrimination, career interruptions from child-bearing, gendered sorting into formal and informal employment, and differences in care responsibilities. By expanding the skill-type space from two types ($\theta_L$, $\theta_H$) to four ($\theta_L^M$, $\theta_H^M$, $\theta_L^F$, $\theta_H^F$), calibrated with sex-specific earnings trajectories and the survival differentials documented in our data section---where women outlive men by 6--7 years across all three countries---the model can decompose fiscal sustainability by sex within the existing equilibrium framework. This decomposition asks: what fraction of the pension system's fiscal pressure comes from women's longer survival (drawing benefits for more periods) versus their lower contribution base (reflecting lower lifetime earnings)? The health expenditure decomposition operates analogously---women's later health crossing point but longer post-retirement survival generates a fiscal profile distinct from men's earlier deterioration but shorter lifespan. This calibration-based approach to gender lets the data carry the information that a structural bargaining model would generate endogenously, at a fraction of the computational cost and with direct empirical grounding.

The paper is structured as follows. Section~\ref{sec:lit} reviews the relevant literature on health financing, OLG models of aging, the longevity channel, formalization incentives, and gender heterogeneity in lifecycle models. Section~\ref{sec:model} presents the theoretical framework---a stochastic OLG economy with endogenous health investment, heterogeneous agents, and a PAYG pension system. Section~\ref{sec:calibration} describes the calibration strategy. Section~\ref{sec:data} documents the data sources for the three country cases. Section~\ref{sec:policy} develops the policy counterfactual framework, and Section~\ref{sec:formalization} presents the formalization experiment. Section~\ref{sec:conclusion} concludes by summarizing the model's analytical contributions and presenting the research plan for the remaining phases of the project.

%% ====================================================================
%% SECTION 2 — LITERATURE REVIEW
%% ====================================================================
\section{Literature Review}
\label{sec:lit}

This paper connects five strands of literature: the economics of health financing and social insurance, the quantitative OLG literature on health and aging, the emerging work on endogenous health investment and longevity as macroeconomic channels, the empirical literature on informality and formalization incentives in Latin America, and the nascent efforts to incorporate gender heterogeneity into lifecycle fiscal models.

%% --------------------------------------------------------------------
\subsection{Health Financing, Social Insurance, and the Tax Wedge}
\label{subsec:lit_health}

The economics of health insurance in developing economies builds on the foundational insight that health shocks are both large and poorly insurable through private markets, generating welfare losses that can be substantially mitigated by pooling arrangements \citep{Arrow1963}. In Latin America, the design question is not whether to pool health risk but through what institutional channel. Systems organized around contributory social security---the Bismarckian tradition---link coverage to formal employment; systems organized around general taxation---the Beveridgean tradition---link coverage to residency or citizenship. In practice, as emphasized by \citet{Levy2008}, most countries in the region operate hybrid systems in which both channels coexist, often with fragmentation, duplication, and unequal benefit packages.

\citet{Summers1989} showed that the welfare implications of mandated benefits depend critically on whether workers value those benefits enough to accept compensating wage reductions. When they do, mandated contributions are essentially priced into wages and generate no net distortion; when they do not---either because benefits are not credible or because workers discount future benefits heavily---contributions function as a tax on formal labor, raising the cost of compliance and reducing formalization. \citet{WorldBankInformalizing2025} formalize this logic through the concept of the Formalization Tax Rate (FTR), which captures the net fiscal cost of formal employment including both mandatory contributions and foregone non-contributory benefits. Their simulations for a range of LAC countries demonstrate that the FTR is highest for low-income workers, precisely the group with the weakest attachment to the formal sector. The FTR framework provides the conceptual bridge between our model's tax wedge---the gap between gross labor cost and net take-home pay for formal versus informal workers---and the observed persistence of informality across the earnings distribution.

The empirical measurement of the tax wedge is provided by the OECD's \textit{Taxing Wages} series. \citet{OECDTaxingWages2026} report that in 2025 the total tax wedge for a single worker earning the average wage in Mexico was 21.7 percent of labour costs---composed of 11.8 percent income tax, 1.4 percent employee social security contributions, and approximately 8.5 percent employer contributions---compared with an OECD average of 35.1 percent. Costa Rica's tax wedge, at 9.8 percent, is among the lowest in the OECD. Two features of the Mexican data are particularly relevant for our calibration. First, the 2021 Social Security Law reform introduced a graduated employer contribution schedule for discharge and old-age insurance (\textit{Cesant\'ia en Edad Avanzada y Vejez}) that scales from 3.15 percent at one minimum wage to 11.875 percent at 4.5 UMAs by 2030, creating a progressive structure in which the marginal cost of formalizing a worker depends on the wage level. Second, the OECD's new progressivity indicator---which measures how the marginal tax wedge varies from 50 to 500 percent of the average wage---places Mexico and Costa Rica at the bottom of the OECD ranking, confirming that the tax system's interaction with informality operates not only through the level of the wedge but through its incidence across the wage distribution.

A central insight from this literature, and one that directly motivates our model's structure, is the asymmetry between health and pension benefits in workers' valuations. Health coverage generates immediate, tangible protection against catastrophic expenditure shocks; pension benefits, by contrast, deliver returns far in the future and are subject to credibility discounts that grow with institutional opacity and long vesting periods. \citet{BaezaPackard2006} argue that health coverage is the social insurance benefit that workers learn most about and value most highly---making health financing design a first-order determinant of compliance behavior. This asymmetry implies that the institutional channel through which health coverage is provided---contributory or tax-financed---has direct consequences for formalization incentives that are larger, per unit of fiscal cost, than those of pension design.

%% --------------------------------------------------------------------
\subsection{OLG Models of Fiscal Sustainability, Health, and Aging}
\label{subsec:lit_olg}

The quantitative OLG literature provides the methodological foundation for this paper. Since \citet{AuerbachKotlikoff1987}, the workhorse framework for analyzing the intergenerational fiscal effects of demographic change has been the multi-period OLG model with a government sector. The health-augmented extensions most relevant to this paper include the work of \citet{DeNardi2010}, who show that uncertain medical expenses are a primary driver of late-life savings among older households in the United States, and \citet{KuhnFeichtinger2015}, who analyze optimal health and retirement choices over the life cycle in a framework where health capital depreciates with age and can be partially restored through investment.

The Grossman tradition of modeling health as a depreciable capital stock---in which individuals invest to maintain or restore a stock that deteriorates with age---provides the theoretical foundation for our health formation equation. \citet{DalgaardStrulik2014} extend this framework to show how optimal aging and mortality patterns emerge endogenously from the interaction of health investment decisions with biological depreciation, yielding a structural interpretation of the Preston curve. \citet{Fioroni2010} analyzes optimal savings and health spending jointly over the life cycle, demonstrating that the interaction between precautionary savings motives and health investment generates consumption-savings profiles that differ substantially from models without endogenous health. \citet{FrankovicWrzaczek2020} study how medical progress interacts with overlapping-generations dynamics, showing that improvements in health technology alter not only individual life-cycle profiles but also aggregate capital accumulation and factor prices through general-equilibrium channels. \citet{YewZhang2018} introduce longevity externalities into a lifecycle-dynastic model, demonstrating that health spending decisions by one generation affect the survival and savings behavior of subsequent generations---a mechanism that, while not directly present in our stationary-equilibrium framework, underscores the importance of modeling health as an endogenous state variable rather than an exogenous process.

For the Latin American context, \citet{CortesEtAl2024} develop a partial-equilibrium OLG model for M\'exico in which agents differ by skill and face stochastic health transitions, health-dependent productivity, and age-varying out-of-pocket expenditure profiles. Their results---consumption gaps of 40--70 percent between good and poor health groups, and a 20-year gap in healthy longevity by skill level---form the empirical backbone of the present paper. The primary extension in this work is the move to general equilibrium with endogenous wages, an explicit pension system with endogenous contribution rates, and a continuous health capital stock with diminishing-returns investment technology. \citet{Lim2020} analyzes the case of South Korea, showing that public provision of health insurance alters aggregate saving behavior in an OLG framework with endogenous health risk---a result that resonates with our model's mechanism, in which the scope of public health coverage affects both household savings decisions and the fiscal cost of the pension system through the contribution base.

%% --------------------------------------------------------------------
\subsection{Endogenous Health Investment and the Longevity Channel}
\label{subsec:lit_longevity}

A recent and rapidly developing strand of the literature studies health investment not merely as a household choice but as a macroeconomic transmission channel. \citet{Carrasco2026} provides a particularly relevant contribution for our framework. In his model, households' health investments endogenously determine longevity, which in turn affects aggregate savings, capital accumulation, and the effectiveness of macroeconomic policy. The key insight is that health investment creates a \textit{longevity channel}: when households invest more in health, they live longer, which strengthens the life-cycle savings motive, raises the aggregate capital stock, and alters the equilibrium interest rate. This channel operates in our model through the interaction of the health formation equation---where health investment $m_j$ augments the stock $h$ that governs both survival probabilities and labor productivity---and the capital market clearing condition, which equates domestic household savings to the aggregate capital stock plus government debt.

The relevance of the longevity channel for Latin America is amplified by two structural features. First, the premature health crossing point documented in Section~\ref{sec:intro}---the fact that health deterioration begins a decade earlier than in high-income countries---implies that the longevity channel operates asymmetrically across skill groups: low-skilled workers, who face earlier and steeper health depreciation, have shorter effective planning horizons and accumulate less precautionary savings, reinforcing the health-income poverty trap that our simulations identify. Second, the interaction between endogenous longevity and the PAYG pension system creates a fiscal amplification mechanism: healthier workers who live longer draw pension benefits for more periods, raising the system dependency ratio $N^R/N^W$ and the endogenous contribution rate $\tau^p$, which in turn increases the tax wedge on formal labor income. This feedback loop---from health investment to longevity to pension costs to the fiscal wedge---is the channel through which spending smarter on health can yield fiscal dividends beyond the direct reduction in out-of-pocket expenditures.

%% --------------------------------------------------------------------
\subsection{Informality, Formalization Incentives, and Social Security}
\label{subsec:lit_informality}

The empirical literature on formalization incentives in Latin America documents a robust negative relationship between social security contribution rates and formal employment, particularly at the lower end of the wage distribution. \citet{Kugler2017} estimate payroll tax elasticities in Colombia ranging between 0.4 and 1.0; similar magnitudes have been found in studies of Turkey \citep{AcarTansel2019}, Brazil, and Mexico. The perceived actuarial value of pension contributions---not just their cost---shapes formalization decisions: evidence from Colombia shows that an increase in the vesting period required to qualify for the minimum pension reduced formal employment by approximately 3 percentage points, illustrating that workers respond to the expected return on their contributions, not merely to the contribution rate itself.\footnote{This result is consistent with the ``social security wealth'' concept developed by \citet{WorldBankInformalizing2025}, which measures the actuarial present value of the benefit stream a worker expects to receive conditional on maintaining formal-sector contributions. Where social security wealth is low---because of long vesting periods, benefit opacity, or low institutional trust---the FTR rises and formalization incentives weaken.}

The non-contributory benefit channel operates in the opposite direction: policies that make health or income transfers available outside formal employment raise the opportunity cost of formalization and tend to shift labor toward the informal sector. \citet{BoschCampos2014} and \citet{AteridoEtAl2011} document significant informality effects from Mexico's Seguro Popular program, though results are not uniform---\citet{AzuaraMarinescu2013} find smaller or negligible effects, suggesting that institutional context and program design mediate the relationship. \citet{CamachoConoverHoyos2013} find analogous results for Colombia's subsidized health insurance expansion, with informality increasing by approximately 4 percentage points. \citet{CalderonMarinescu2012} present a more nuanced picture for Colombia, finding that regulations unifying health and pension payment systems partially offset the informality-inducing effects of coverage expansion. The general lesson, consistent with the analytical framework of this paper, is that the design of the health benefit package---in particular the degree to which it is tied to formal contributions versus financed through general taxation---is a first-order determinant of formalization equilibria.

The broader structural context is provided by \citet{PerryEtAl2007}, who argue that widespread informality reflects a mass opting out of state institutions---a ``blunt societal indictment of the quality of the State's services provision.'' \citet{WorldBankInformalizing2025} extend this line of analysis to its most comprehensive treatment to date, arguing that three dimensions must move together for formalization to improve: market competitiveness (the tax wedge must not exceed what firms can absorb), workers' valuation of contributory benefits (``social security wealth'' must be perceived as positive), and enforcement credibility (compliance must be monitored and sanctioned). Their reform roadmap emphasizes that LAC countries face a compressed timeline---the aging that took Europe 70 years will take Latin America approximately 25---leaving far less room for incremental adjustment. The report's conclusion that the ``crisis is not yet averted'' provides the institutional backdrop against which our model's policy counterfactuals should be read. \citet{BoschMelguizoPages2013} document how pension design interacts with labor market outcomes across the region, showing that systems with tighter benefit-contribution links generate stronger formalization incentives.

Our model contributes to this literature by calibrating the baseline to the formal-sector population and studying informality through a structured policy experiment: a formalization program that incorporates a fraction of the informal population into the contributory system, using observed income and health profiles from national household surveys. This design allows us to quantify the competing fiscal effects of a broader contribution base against higher per-capita health expenditure claims from newly formalized low-skilled workers, within the general-equilibrium discipline of endogenous factor prices and pension clearing.

%% --------------------------------------------------------------------
\subsection{Gender Heterogeneity in Lifecycle Fiscal Models}
\label{subsec:lit_gender}

The incorporation of gender into quantitative OLG models remains at an early stage. \citet{FehrKindermann2018} provide the most widely used pedagogical framework, constructing a two-person household in which the female member's labor force participation decision is modeled as a function of human capital accumulation, care responsibilities, and the relative returns to market work versus home production. Their model demonstrates the basic mechanism---female participation responds to policy instruments such as childcare subsidies, parental leave, and joint versus individual taxation---but operates in a substantially simpler economic environment than the one developed in this paper: health is not a state variable, there is no endogenous health investment or survival risk, no longevity channel, and no formalization margin.

The computational obstacle to combining household bargaining with rich health dynamics is severe. \citet{DeNardi2010} and subsequent work in the same line have extended models of medical expense risk to two-person households---singles and couples---but in a partial-equilibrium setting, with factor prices treated as exogenous. Embedding such a two-person household structure in a general-equilibrium framework, where wages and the interest rate respond endogenously to aggregate savings and labor supply, raises the computational burden substantially: the curse of dimensionality compounds with the fixed-point problem in factor prices. With the continuous health capital, age-dependent depreciation, endogenous survival, and endogenous factor prices that characterize our model, a structural household bargaining extension would be computationally challenging and, to our knowledge, has not been attempted in the published literature.

Our approach to gender is different. Rather than modeling the household decision-making process, we exploit the fact that sex-differentiated income profiles and survival probabilities already embed---in reduced form---the cumulative effects of labor market discrimination, career interruptions from child-bearing, gendered sorting into formal and informal employment, and the unequal distribution of care responsibilities. By expanding the agent-type space from two skill types $(\theta_L, \theta_H)$ to four $(\theta_L^M, \theta_H^M, \theta_L^F, \theta_H^F)$, calibrated with sex-specific earnings trajectories from household surveys and the survival differentials documented in national life tables, the model can decompose fiscal sustainability pressures by sex within the existing equilibrium framework. This calibration-based strategy lets the data carry the information that a structural bargaining model would generate endogenously, at a fraction of the computational cost and with direct empirical grounding. The structural household model remains a natural extension for future research, but it is not a prerequisite for the first-order fiscal decomposition that this paper's policy questions require.

\bigskip
%% --------------------------------------------------------------------

\section{Model}
\label{sec:model}

\bigskip

Section~\ref{sec:model} presents the theoretical framework that integrates these five strands. The model is a stochastic OLG economy with endogenous health investment, heterogeneous agents, and an explicit PAYG pension system; it extends the ITED baseline \citep{CortesEtAl2024} to incorporate general equilibrium factor pricing, a continuous health capital stock with age-dependent depreciation, and a contribution-rate rule that endogenously links demographic aging to the fiscal wedge on formal labor.
\subsection{Environment and Demographics}
\label{subsec:environment}

Time is discrete. The economy is populated by $J$ overlapping generations indexed by $j = 1, \ldots, J$, where $j = 1$ corresponds to age~20 and each model period spans five years. Households enter economic life at $j = 1$, retire exogenously at age $j_R$, and may live until $j = J$ (age~100).

\paragraph{Agent types.}
The economy contains households of two permanent skill types, $\theta \in \{\theta_L, \theta_H\}$, where $\theta_L$ denotes low-skilled and $\theta_H$ high-skilled. Skill type is drawn at birth and remains fixed over the life cycle. It governs labor productivity through the exposure parameter $\varrho(\theta)$ in the productivity function~\eqref{eq:labor_productivity}. There is a unit mass of newborns of each type each period and the population grows at gross rate $(1 + n_p)$.

\paragraph{Idiosyncratic risk.}
Households face two sources of individual uncertainty: (i)~mortality risk, depending on health, and (ii)~idiosyncratic labor-productivity risk, summarized by an AR(1) shock $\eta$. Health $h$ is endogenous: it evolves deterministically given current health and chosen medical spending. There is no aggregate uncertainty.

\paragraph{Individual state.}
The individual state of a household at the beginning of a period is
\begin{equation}
\label{eq:state}
s = (a,\, h,\, \eta),
\end{equation}
where $a \in \mathcal{A} = [0, \bar{a}]$ are beginning-of-period assets, $h \in \mathcal{H}_h = [0, \bar{h}]$ is health capital, and $\eta \in \mathcal{H}_\eta = \{\eta_1, \ldots, \eta_{N_\eta}\}$ is the idiosyncratic productivity component (discretized from an AR(1) process; see Section~\ref{subsubsec:productivity}). Skill type $\theta$ and age $j$ are deterministic indices that label the value and policy functions.

\paragraph{Cohort weights.}
Let $\phi_j(\theta)$ denote the stationary measure of age-$j$ households of type~$\theta$, normalized so that $\sum_\theta \sum_{j=1}^{J} \phi_j(\theta) = 1$ and $\phi_1(\theta) = \phi_1$ for both types (equal birth shares). These weights satisfy
\begin{equation}
\label{eq:cohort_weights}
\phi_{j+1}(\theta) \;=\; \frac{\bar{\psi}_{j+1}(\theta)}{1 + n_p}\, \phi_j(\theta), \qquad j = 1, \ldots, J-1,
\end{equation}
% CORRECTION: The original definition averaged survival over an EXOGENOUS
% Markov kernel \pi_j^m for health. But health is endogenous in this model,
% so the average must be taken over the equilibrium distribution. The correct
% formula is:
where the population-weighted average survival is
\begin{equation}
\label{eq:avg_survival}
\bar{\psi}_{j+1}(\theta)
\;=\;
\frac{1}{\phi_j(\theta)} \int \psi_{j+1}\bigl(h'(s; j, \theta)\bigr)\, d\tilde\mu_j^\theta(s),
\end{equation}
with $h'(s; j, \theta) = (1-\delta_j^h)h + H_j(m(s; j, \theta))$ the endogenous next-period health and $\tilde\mu_j^\theta$ the unconditional age-$j$ population mass for type $\theta$ (so that $\int d\tilde\mu_j^\theta = \phi_j(\theta)$). $\psi_J(h) \equiv 0$ for all $h$ (certain death at the terminal age).
% CORRECTION: Made the integrating measure unconditional so the formula is
% well-defined (and consistent with the law of motion below, where the
% measure mu carries cohort mass directly rather than being a conditional
% probability).
 

\paragraph{Within-period timing.}
The timing within each period is:
\begin{enumerate}
\item The household observes its current state $(a, h, \eta)$.
\item Working-age households choose $(c, \ell, a', m)$. Retired households choose $(c, a', m)$.
\item End-of-period assets are $a'$. Survival to next period is realized with probability $\psi_{j+1}(h')$, where $h' = (1-\delta_j^h) h + H_j(m)$ is determined by the chosen $m$. We assume that assets of agents who die are voided (do not enter aggregate capital next period; do not flow to survivors).
\end{enumerate}


\subsection{Preferences}
\label{subsec:preferences}

Households rank streams of consumption, leisure, and health by expected discounted lifetime utility. The lifetime objective of a newborn household of type~$\theta$ is
\begin{equation}
\label{eq:lifetime_utility}
\E_1 \left[\sum_{j=1}^{J} \beta^{j-1} \left(\prod_{i=2}^{j} \psi_i(h_{i-1}')\right) u(c_j,\, \ell_j,\, h_j)\right],
\end{equation}
where $\beta \in (0,1)$, the survival probability $\psi_i(h_{i-1}')$ is the probability of being alive at age $i$ given that next-period health entering age $i$ is $h_{i-1}' = h_i$, and the expectation is taken over the path of $\eta$ realizations.
 
The period utility function is GHH-augmented with a health amenity:
\begin{equation}
\label{eq:utility}
u(c, \ell, h) \;\equiv\; \frac{\bigl(c + s(h) - v(\ell)\bigr)^{1-\gamma} - 1}{1 - \gamma},
\qquad
v(\ell) \equiv \Psi \frac{\ell^{1+\nu}}{1+\nu},
\end{equation}
\begin{equation}
    s(h) \equiv
\begin{cases}
\Xi \dfrac{h^{1-\xi}}{1-\xi} & \xi \neq 1,\\[6pt]
\Xi \log(h) & \xi = 1,
\end{cases}
\end{equation}

with $\Psi > 0$, $\Xi \geq 0$, $\nu > 0$, $\gamma > 0$. Here $c \geq 0$ is consumption, $\ell \in [0, 1]$ is hours worked ($1-\ell$ is leisure), and $h$ is the contemporaneous stock of health capital that simultaneously affects survival, productivity, and the amenity value of consumption.

\begin{remark}[Retirees]
\label{rem:retirees}
For retirees ($j \geq j_R$), $\ell_j = 0$ and the period utility reduces to $u(c, 0; h) = \frac{(c + s(h))^{1-\gamma} - 1}{1-\gamma}$.
\end{remark}


\subsection{Technology}
\label{subsec:technology}

A representative, competitive firm produces the single consumption good using capital $K$ and aggregate effective labor $L$ according to the Cobb-Douglas technology
\begin{equation}
\label{eq:production}
Y = A\, K^\alpha\, L^{1-\alpha}, \qquad \alpha \in (0,1),
\end{equation}
with $A > 0$ total factor productivity. Capital depreciates at rate $\delta \in (0,1)$. Because the production function has constant returns to scale, the firm's first-order conditions depend only on the ratio $K/L$ and so are invariant to the choice of units. We will write $K$ and $L$ in stationary per-cohort-mass units throughout (see notation block in Section~\ref{subsec:markets}); this convention is consistent with the firm's optimization.
 
\paragraph{Firm's problem.}
The firm takes factor prices $(r, w)$ as given and maximizes static profits:
\begin{equation}
\label{eq:firm_problem}
\max_{K, L} \bigl\{A K^\alpha L^{1-\alpha} - (r+\delta)\,K - w\,L\bigr\}.
\end{equation}
The first-order conditions yield
\begin{align}
r &= \alpha A \left(\frac{K}{L}\right)^{\alpha-1} - \delta, \label{eq:foc_r} \\
w &= (1-\alpha)\, A \left(\frac{K}{L}\right)^{\alpha}. \label{eq:foc_w}
\end{align}
where $\delta \in (0,1)$ is the rate of capital depreciation.

\paragraph{Factor prices.}
Both $r$ and $w$ are determined endogenously in equilibrium. Given aggregate capital $K$ and aggregate labor $L$, equations~\eqref{eq:foc_r} and \eqref{eq:foc_w} determine factor prices. Equivalently, the capital-to-labor ratio $K/L$ pins both prices simultaneously:
\begin{align}
r &= \alpha A \left(\frac{K}{L}\right)^{\alpha - 1} - \delta, \label{eq:interest_rate} \\
w &= (1-\alpha)\, A \left(\frac{K}{L}\right)^{\alpha}. \label{eq:wage}
\end{align}
The equilibrium interest rate is the object that adjusts to clear the capital market.


\subsection{Exogenous Processes}
\label{subsec:exogenous}

\subsubsection{Mortality}
\label{subsubsec:mortality}

Conditional survival probabilities depend on age and current health. The probability of surviving from age $j-1$ to age $j$, known at the start of age $j-1$, is $\psi_j(m_{j-1})$, with $\psi_J(m) = 0$ for all $m$ (certain death at the terminal age).

\subsubsection{Labor Productivity}
\label{subsubsec:productivity}

Effective labor productivity of a working-age household ($j < j_R$) is
\begin{equation}
\label{eq:labor_productivity}
\nu_j(h, \eta;\, \theta) \;=\; e_j \exp\bigl(\theta + \eta - \varrho(\theta)\, h\bigr),
\end{equation}
where $e_j > 0$ is a deterministic age-efficiency profile and $\varrho(\theta) \geq 0$ is a health-productivity penalty. Productivity is zero for retirees ($j \geq j_R$).

The idiosyncratic component $\eta$ follows the AR(1) process
\begin{equation}
\label{eq:ar1}
\eta_{j+1} = \rho\, \eta_j + \varepsilon_{j+1}, \qquad \varepsilon_{j+1} \sim \mathcal{N}(0, \sigma_\varepsilon^2),
\end{equation}
with $|\rho| < 1$. 

\subsubsection{Health Formation and Health Investment}
\label{subsubsec:health}

Health evolves deterministically given the chosen medical spending:
\begin{equation}
\label{eq:health_formation}
h' = \min\bigl\{(1 - \delta_j^h)\, h + H_j(m),\; \bar h\bigr\},
\end{equation}
where the depreciation rate $\delta_j^h \in (0, 1)$ is age-dependent and $H_j: \mathbb{R}_+ \to \mathbb{R}_+$ satisfies $H_j(0) = 0$, $H_j' > 0$, $H_j'' < 0$ (diminishing returns to medical spending). A specific functional form (e.g., $H_j(m) = \xi_j m^{\zeta}$ with $\zeta \in (0,1)$) must be committed to before computation.


\subsection{Government}
\label{subsec:government}

The government provides three programs: general fiscal policy (taxes, public goods, and a consumption floor), and a pay-as-you-go (PAYG) pension system. We describe each in turn.

\subsubsection{Fiscal Policy}
\label{subsubsec:fiscal}
The government taxes consumption at rate $\tau^c$, labor income at rate $\tau^\omega$, capital income at rate $\tau^k$, and medical spending at rate $\tau^m$. It purchases public goods $G = g_y Y$ and may issue one-period debt $B$. All aggregate quantities ($C$, $K$, $L$, $M$, $Y$, $G$, $B$) are in stationary per-cohort-mass units (see notation block in Section~\ref{subsec:markets}); $C$, $L$, $M$, $Y$, $G$ are per-period flows and $K$, $B$ are stocks. The flow budget is:
\begin{equation}
\label{eq:gov_budget}
\tau^c C + \tau^m M + \tau^\omega w L + \tau^k r K + B' \;=\; G + (1+r)B,
\end{equation}
where $C$ is aggregate consumption, $K$ is the aggregate capital stock, $M$ is aggregate health investments, and $B'$ denotes next-period government debt.

\paragraph{Fiscal closure.}
$(\tau^c, \tau^\omega, \tau^k, \tau^m, g_y)$ are exogenous policy parameters. In the stationary equilibrium $B' = B$, so debt service $r B$ equals the primary surplus:
\begin{equation}
\label{eq:gov_stationary}
r\, B \;=\; \tau^c\, C + \tau^\omega\, w\, L + \tau^k\, r\, K + \tau^m\, M - G.
\end{equation}
$B$ is the residual that closes the budget. (In counterfactual exercises, the user must verify whether $B$ is the appropriate clearing instrument; alternatives include adjusting one tax rate or $g_y$.)

\subsubsection{Pension System}
\label{subsubsec:pension}
A PAYG pension is financed by a payroll tax $\tau^p$ on workers' labor income. The system balances each period:
\begin{equation}
\label{eq:pension_balance}
\tau^p\, w\, L \;=\; \overline{pen}\, N^R,
\end{equation}
where $N^R = \sum_\theta \sum_{j=j_R}^{J} \phi_j(\theta)$ is the mass of retirees. The pension benefit is the replacement-rate parameter $\kappa$ times average per-worker labor income:
\begin{equation}
\label{eq:pension_benefit}
\overline{pen} \;=\; \kappa \cdot \frac{w L}{N^W},
\end{equation}
where $N^W = \sum_\theta \sum_{j=1}^{j_R - 1} \phi_j(\theta)$ is the mass of working-age individuals.

\paragraph{Pension closure.}
The replacement rate $\kappa$ is an exogenous policy parameter. The contribution rate $\tau^p$ adjusts endogenously to balance the pension system each period. Substituting \eqref{eq:pension_benefit} into \eqref{eq:pension_balance} yields
\begin{equation}
\label{eq:pension_contribution}
\tau^p \;=\; \kappa\, \frac{N^R}{N^L},
\end{equation}
so the contribution rate equals the replacement rate scaled by the ratio of pension-eligible retirees to formal workers (the system dependency ratio). This relationship pins down $\tau^p$ once the stationary distribution is known.


\subsection{Markets}
\label{subsec:markets}
\par Throughout this section we work in \emph{stationary per-cohort-mass} units. Specifically: aggregates are integrals of individual quantities against the stationary population measure $\mu_j^\theta$, with $\sum_\theta \sum_j \phi_j(\theta) = 1$. Population growth $n_p$ is already absorbed into the cohort-weight recursion~\eqref{eq:cohort_weights}; the aggregates defined below therefore do \emph{not} grow with the actual population. A symbol such as $K$ in the Markets and Equilibrium sections refers to the stationary capital stock per unit of cohort mass, related to the absolute capital stock $K^{\mathrm{abs}}_t$ at calendar time $t$ by $K^{\mathrm{abs}}_t = K \cdot N_t$, where $N_t$ is total population at $t$ growing at gross rate $1+n_p$. We further distinguish stocks from per-period flows. \emph{Stocks} (units: goods, or units of agent-mass): $K$, $A_{\mathrm{dom}}$, $B$, $N^W$, $N^R$. \emph{Per-period flows} (units: goods per period, or efficiency units of labor per period): $Y$, $C$, $M$, $G$, $\delta K$, $L$, $\Lambda_{\mathrm{void}}$, $\overline{pen}$. The goods-market identity~\eqref{eq:goods_market} is a flow identity: every term has units of goods per period.

\par Asset markets are incomplete. Households can save in a single non-contingent asset $a$ and are not allowed to borrow, so
\begin{equation}
\label{eq:borrowing_constraint}
a' \;\geq\; 0.
\end{equation}
Both factor prices $(r, w)$ are determined endogenously by equations~\eqref{eq:interest_rate}--\eqref{eq:wage}. Markets clear when the following conditions hold:

\paragraph{Labor markets.}
\begin{equation}
\label{eq:labor_market}
L \;=\; \sum_\theta\, \sum_{j=1}^{j_R - 1}\, \int \nu_j(h, \eta;\, \theta)\, \ell(s;\, j, \theta)\, d\mu_j^\theta(s),
\end{equation}
\paragraph{Capital market.}
Domestic household savings finance the entire capital stock and government debt:
\begin{equation}
\label{eq:capital_market}
K \;=\; A_{\mathrm{dom}} - B,
\qquad
A_{\mathrm{dom}} \;=\; \sum_\theta\, \sum_{j=1}^{J}\, \int a\, d\mu_j^\theta(s).
\end{equation}
This condition, together with the firm's FOC \eqref{eq:interest_rate}, determines the equilibrium interest rate.

\paragraph{Goods market.}
\begin{equation}
\label{eq:goods_market}
Y \;=\; C + M + \delta\, K + G + \Lambda_{\mathrm{void}},
\end{equation}
where every term is a flow (goods per period in stationary per-cohort-mass units), and the voided-asset flow is
\begin{equation}
\label{eq:leakage_voided}
\Lambda_{\mathrm{void}}
\;\equiv\;
\sum_\theta\, \sum_{j=1}^{J}\, \bigl(1 - \bar\psi_{j+1}(\theta)\bigr)
\int a'(s;\, j, \theta)\, d\mu_j^\theta(s).
\end{equation}
\paragraph{Government bond market.}
By construction the government issues exactly $B$ bonds; households' aggregate asset holdings finance both $K$ and $B$:
\begin{equation}
A_{\mathrm{dom}} = K + B,
\end{equation}
which is~\eqref{eq:capital_market} rearranged.

\subsection{Optimization Problems}
\label{subsec:optimization}

\subsubsection{Net Labor Income}
\label{subsubsec:net_income}

Define net labor income as
\begin{equation}
\label{eq:net_labor_income}
y(h, \eta, \ell;\, \theta) \;=\; w\, \nu(h, \eta;\, \theta)\, \ell\, \bigl(1 - \tau^\omega - \tau^p\bigr).
\end{equation}
Formal workers pay labor income tax $\tau^\omega$ and pension contributions $\tau^p$; informal workers pay neither.

\subsubsection{Workers' Problem}
\label{subsubsec:worker}

Individuals observe their productivity shock and then choose consumption, savings, labor supply, and health investment. Given labor supply $\ell$, define household resources as
\begin{equation}
\label{eq:resources_worker}
x(a, h, \eta, \ell;\, \theta) \;=\; \bigl(1 + r(1 - \tau^k)\bigr) a + y(h, \eta, \ell;\, \theta).
\end{equation}
The worker's Bellman equation is
\begin{equation}
\label{eq:worker_bellman}
V_j(a, h, \eta;\, \theta)
\;=\;
\max_{\substack{c \geq 0,\, a' \geq 0 \\ \ell \in [0,1],\, m \geq 0}}
\Bigl\{
u(c, \ell, h)
+ \beta\, \psi_{j+1}(h')\,
\E\bigl[V_{j+1}(a', h', \eta';\, \theta) \mid \eta\bigr]
\Bigr\}
\end{equation}
subject to
\begin{align}
(1 + \tau^c)\, c + (1 + \tau^m)\, m + a' &= x(a, h, \eta, \ell;\, \theta), \label{eq:worker_budget}\\
h' &= \min\bigl\{(1 - \delta_j^h)\, h + H_j(m),\, \bar h\bigr\}, \label{eq:worker_health}\\
a' &\geq 0,\quad m \geq 0,\quad \ell \in [0,1]. \label{eq:worker_borrow}
\end{align}
Let the optimal savings policy be denoted by $a_j'(a, h, \eta;\, \theta)$.

\subsubsection{Retiree's Problem}
\label{subsubsec:retiree}
\begin{equation}
\label{eq:resources_retiree}
x(a; \theta) \;=\; \bigl(1 + r(1-\tau^k)\bigr) a + \overline{pen}.
\end{equation}
The Bellman equation is
\begin{equation}
\label{eq:retiree_bellman}
V_j(a, h, \eta;\, \theta)
\;=\;
\max_{c \geq 0,\, a' \geq 0,\, m \geq 0}
\Bigl\{
u(c, 0, h)
+ \beta\, \psi_{j+1}(h')\,
\E\bigl[V_{j+1}(a', h', \eta';\, \theta) \mid \eta\bigr]
\Bigr\}
\end{equation}
subject to
\begin{align}
(1 + \tau^c)\, c + (1 + \tau^m)\, m + a' &= x(a; \theta), \label{eq:retiree_budget}\\
h' &= \min\bigl\{(1 - \delta_j^h)\, h + H_j(m),\, \bar h\bigr\}, \label{eq:retiree_health}\\
a' &\geq 0,\quad m \geq 0. \label{eq:retiree_borrow}
\end{align}
At the terminal age $j = J$, $\psi_{J+1}(h') = 0$, so the continuation value vanishes and the optimum is $a'_J = 0$, $m_J = 0$ (no point investing in health when survival probability is zero).

\subsubsection{First-Order Conditions}
\label{subsubsec:focs}

Let $\lambda_j$ denote the multiplier on the budget constraint. The interior FOCs are:
\begin{align}
\text{Consumption:} \quad & u_c(c, \ell, h) \;=\; \lambda_j (1 + \tau^c), \label{eq:foc_c}\\
\text{Labor:} \quad & -u_\ell(c, \ell, h) \;=\; \lambda_j\, w\, \nu_j(h, \eta; \theta)(1-\tau^\omega-\tau^p), \label{eq:foc_l}\\
\text{Savings:} \quad & \lambda_j \;=\; \beta\, \psi_{j+1}(h')\, \E\bigl[V_{a, j+1}(a', h', \eta'; \theta) \mid \eta\bigr], \label{eq:foc_a}\\
\text{Medical spending:} \quad & \lambda_j (1 + \tau^m)
\;=\; \beta\, H_j'(m)\,\Bigl\{
\psi_{j+1}'(h')\, \E[V_{j+1}(a', h', \eta'; \theta) \mid \eta] \notag\\
&\qquad\qquad\qquad\qquad +\, \psi_{j+1}(h')\, \E[V_{h, j+1}(a', h', \eta'; \theta) \mid \eta]
\Bigr\}. \label{eq:foc_m}
\end{align}

\paragraph{Envelope conditions.} For the worker's problem,
\begin{align}
V_{a, j}(a, h, \eta; \theta) &\;=\; \lambda_j \cdot \bigl(1 + r(1-\tau^k)\bigr), \label{eq:env_a}\\
V_{h, j}(a, h, \eta; \theta) &\;=\; u_h(c, \ell, h) + \lambda_j\, w\, \nu_{j,h}(h, \eta; \theta)\, \ell\, (1-\tau^\omega-\tau^p) \notag\\
&\qquad +\, \beta\, (1-\delta_j^h)\, \psi_{j+1}'(h')\, \E[V_{j+1}(a', h', \eta'; \theta) \mid \eta] \notag\\
&\qquad +\, \beta\, (1-\delta_j^h)\, \psi_{j+1}(h')\, \E[V_{h, j+1}(a', h', \eta'; \theta) \mid \eta]. \label{eq:env_h}
\end{align}

\paragraph{Euler equation.} Combining \eqref{eq:foc_a} with \eqref{eq:env_a} (one period ahead) and \eqref{eq:foc_c}:
\begin{equation}
\label{eq:euler}
u_c(c, \ell, h) \;=\; \beta\, \bigl(1 + r(1-\tau^k)\bigr)\, \psi_{j+1}(h')\, \E\bigl[u_c(c', \ell', h') \mid \eta\bigr].
\end{equation}
 
\paragraph{Consumption--leisure condition.} Combining \eqref{eq:foc_c} and \eqref{eq:foc_l}:
\begin{equation}
\label{eq:cons_leisure}
\frac{-u_\ell(c, \ell, h)}{u_c(c, \ell, h)}
\;=\;
\frac{w\, \nu_j(h, \eta; \theta)(1-\tau^\omega-\tau^p)}{1 + \tau^c}.
\end{equation}
 
\paragraph{Kuhn-Tucker for the borrowing constraint.}
\begin{equation}
\bigl(\lambda_j - \beta\, \psi_{j+1}(h')\, \E[V_{a, j+1}(a', h', \eta'; \theta) \mid \eta]\bigr)\, a' = 0,\quad a' \geq 0.
\end{equation}
 


\subsection{Stationary Competitive Equilibrium}
\label{subsec:equilibrium}

Let $\mu_j^\theta$ denote the time-invariant population measure of age-$j$ households of type $\theta$ over the state space $\mathcal{S} = \mathcal{A} \times [0, \bar h] \times \mathcal{H}_\eta$. The convention is that $\mu_j^\theta$ is \textbf{unconditional}: $\int d\mu_j^\theta = \phi_j(\theta)$. Aggregates are integrals over this measure (no extra $\phi_j$ outside).
 
\subsubsection{Aggregate Objects}
 
\begin{align}
C &\;=\; \sum_\theta\, \sum_{j=1}^{J}\, \int c(s; j, \theta)\, d\mu_j^\theta(s), \label{eq:agg_C}\\
M &\;=\; \sum_\theta\, \sum_{j=1}^{J}\, \int m(s; j, \theta)\, d\mu_j^\theta(s), \label{eq:agg_M}\\
L &\;=\; \sum_\theta\, \sum_{j=1}^{j_R - 1}\, \int \nu_j(h, \eta; \theta)\, \ell(s; j, \theta)\, d\mu_j^\theta(s), \label{eq:agg_L}\\
A_{\mathrm{dom}} &\;=\; \sum_\theta\, \sum_{j=1}^{J}\, \int a\, d\mu_j^\theta(s), \label{eq:agg_A}\\
N^W &\;=\; \sum_\theta\, \sum_{j=1}^{j_R - 1}\, \phi_j(\theta), \qquad N^R \;=\; \sum_\theta\, \sum_{j=j_R}^{J}\, \phi_j(\theta), \label{eq:agg_NWR}\\
\Lambda_{\mathrm{void}} &\;=\; \sum_\theta\, \sum_{j=1}^{J}\, \bigl(1 - \bar\psi_{j+1}(\theta)\bigr) \int a'(s; j, \theta)\, d\mu_j^\theta(s). \label{eq:agg_Lvoided}
\end{align}

\subsubsection{Law of Motion for the Distribution}
 
For each measurable set $\mathcal{B} \subseteq \mathcal{S}$ and each age $j = 1, \ldots, J-1$ and type $\theta$:
\begin{equation}
\label{eq:lom}
\mu_{j+1}^\theta(\mathcal{B})
\;=\;
\frac{1}{1 + n_p}
\int
\psi_{j+1}\bigl(h'(s; j, \theta)\bigr)\,
\Bigl[\sum_{\eta' \in \mathcal{H}_\eta}
\Pi^\eta(\eta, \eta')\,
\mathbf{1}\bigl[\bigl(a'(s; j, \theta),\, h'(s; j, \theta),\, \eta'\bigr) \in \mathcal{B}\bigr]
\Bigr]
\, d\mu_j^\theta(s).
\end{equation}
The initial condition for newborns is
\begin{equation}
\label{eq:initial_dist}
\mu_1^\theta(\mathcal{B})
\;=\;
\phi_1(\theta) \cdot
\sum_{\eta_1 \in \mathcal{H}_\eta} \pi^\eta_{\mathrm{erg}}(\eta_1)\,
\mathbf{1}[(0, h_0, \eta_1) \in \mathcal{B}],
\end{equation}
where $\pi^\eta_{\mathrm{erg}}$ is the ergodic distribution of $\Pi^\eta$ and $h_0$ is initial health (a parameter, e.g.\ $h_0 = \bar h$).

\subsubsection{Definition}
 
\begin{definition}[Stationary Recursive Competitive Equilibrium]
\label{def:equilibrium}
Given exogenous policy $(\tau^c, \tau^\omega, \tau^k, \tau^m, g_y, \kappa)$, parameters of preferences, technology, and demographic processes, a \textbf{stationary recursive competitive equilibrium} consists of
\begin{enumerate}
\item value functions $\{V_j(\cdot; \theta)\}_{j = 1, \ldots, J;\, \theta}$ and policy functions $\{c_j(\cdot; \theta), \ell_j(\cdot; \theta), m_j(\cdot; \theta), a'_j(\cdot; \theta)\}$ defined on $\mathcal{S}$;
\item factor prices $(r, w)$;
\item endogenous fiscal/pension variables $(B, \tau^p, \overline{pen})$;
\item aggregates $(K, L, C, M, A_{\mathrm{dom}}, N^W, N^R, \Lambda_{\mathrm{void}})$ in stationary per-cohort-mass units (see notation block in Section~\ref{subsec:markets});
\item cohort weights $\{\phi_j(\theta)\}_{j, \theta}$ and population-average survivals $\{\bar\psi_{j+1}(\theta)\}_{j, \theta}$;
\item stationary distributions $\{\mu_j^\theta\}_{j, \theta}$;
\end{enumerate}
such that:
\begin{enumerate}
\item \textbf{Worker optimality.} For each $\theta$ and $j < j_R$, $V_j$ solves the worker's Bellman~\eqref{eq:worker_bellman}--\eqref{eq:worker_borrow}, with policy functions $(c_j, \ell_j, m_j, a'_j)$.
\item \textbf{Retiree optimality.} For each $\theta$ and $j \geq j_R$, $V_j$ solves the retiree's Bellman~\eqref{eq:retiree_bellman}--\eqref{eq:retiree_borrow}, with policy functions $(c_j, m_j, a'_j)$ and $\ell_j \equiv 0$.
\item \textbf{Firm optimality.} $(r, w)$ satisfy~\eqref{eq:foc_r}--\eqref{eq:foc_w} given $(K, L)$.
\item \textbf{Government budget.} $B$ satisfies~\eqref{eq:gov_stationary}.
\item \textbf{Pension balance.} $\tau^p$ satisfies~\eqref{eq:pension_contribution} and $\overline{pen}$ satisfies~\eqref{eq:pension_benefit}.
\item \textbf{Market clearing.}
\begin{itemize}
\item Labor:~\eqref{eq:labor_market}.
\item Capital:~\eqref{eq:capital_market}.
\item Goods:~\eqref{eq:goods_market} (holds as a Walras-law identity).
\end{itemize}
\item \textbf{Cohort weights.} $\phi_j(\theta)$ satisfy~\eqref{eq:cohort_weights} with $\bar\psi_{j+1}(\theta)$ given by~\eqref{eq:avg_survival}, normalized so that $\sum_\theta \sum_j \phi_j(\theta) = 1$.
\item \textbf{Stationarity of the distribution.} For each $j$ and $\theta$, $\mu_j^\theta$ satisfies the law of motion~\eqref{eq:lom}, with the newborn condition~\eqref{eq:initial_dist}.
\item \textbf{Aggregate consistency.} $(C, M, L, A_{\mathrm{dom}}, N^W, N^R, \Lambda_{\mathrm{void}})$ are computed from~\eqref{eq:agg_C}--\eqref{eq:agg_Lvoided}, all in stationary per-cohort-mass units.
\end{enumerate}
\end{definition}


%% ====================================================================
\section{Calibration Strategy}
\label{sec:calibration}

We calibrate the model in two steps, following \cite{gourinchas2002consumption}. In the first step, we assign values to parameters that can be identified outside the model and estimated directly from the data. In the second step, we internally calibrate the remaining parameters using the Simulated Method of Moments (SMM).

The calibration therefore involves three groups of parameters. The first group consists of parameters taken directly from the literature, including the intertemporal elasticity of substitution and the parameters governing the final-good production function. The second group contains parameters estimated outside the model in the first-step calibration, including government policy parameters and the parameters governing mortality and labor productivity. The third group consists of parameters calibrated internally in the second step, primarily those governing preferences and the health technology.

\subsection{First-Step Calibration}

The first-step calibration estimates the parameters governing the mortality and labor productivity processes.

\subsubsection{Mortality}

We estimate mortality probabilities using a Probit model. Let $d_{i,t}$ denote an indicator equal to one if individual $i$ is dead at time $t$. Conditional on being alive at time $t$, the probability of dying before period $t+1$ is given by
\begin{equation}
\label{eq:probit}
    \Pr(d_{i,t+1}=1 \mid d_{i,t}=0)=\Phi(X_{it}'\gamma),
\end{equation}
where $\Phi(\cdot)$ denotes the standard normal cumulative distribution function. The vector $X_{it}$ includes a second-order polynomial in age, a second-order polynomial in health, and additional demographic controls.

\subsubsection{Labor Productivity}

We approximate individual labor productivity using observed wages, as is standard in the literature. To estimate the deterministic component of the labor productivity process in equation~\eqref{eq:labor_productivity}, we regress log wages on second-order polynomials in age and health, as well as education controls.

We then use the residuals from this regression to estimate the stochastic component of labor productivity. Specifically, we compute the variance-covariance matrix of the residuals and use equally weighted minimum distance to estimate the autoregressive coefficient, $\rho$, and the variance of the innovation to the persistent productivity component, $\sigma_{\epsilon}^{2}$.

\subsection{Second-Step Calibration}

The second-step calibration assigns values to the parameters governing health production and preferences. These parameters are calibrated internally by matching targeted moments in the data.

\subsubsection{Health Production}

For the health production function, we follow \cite{Carrasco2026} and adopt the specification
\begin{equation}
    \label{eq:healthproduction}
    H_{j}(m) = \overline{H}_{j}m^{\zeta_{h}},
\end{equation}
where
\begin{equation}
    \overline{H}_{j}=\overline{H}_{0}\exp(-h^{slope}j).
\end{equation}
The parameter $h^{slope}$ governs the age-related decline in the productivity of health investment. We calibrate the parameters of the health production function to match the life-cycle profile of medical expenditures.

\subsubsection{Preference Parameters}

The period utility function in equation~\eqref{eq:utility} is governed by five preference parameters: the curvature of the GHH composite, $\gamma$; the inverse Frisch elasticity of labor supply, $\nu$; the labor disutility scale, $\Psi$; and the level and curvature of the health amenity, $(\Xi,\xi)$. We discipline these parameters using a combination of external evidence and internal moment matching.

We set the curvature of the GHH composite to $\gamma=2$, a value within the range commonly used in the consumption and asset-pricing literatures. We set the inverse Frisch elasticity to $\nu=2$, which implies a Frisch elasticity of labor supply equal to $0.5$. This value is consistent with the consensus estimate from microeconometric studies of prime-age workers \citep{chetty_etal_2011}. A useful feature of the GHH specification is that $\nu$ is identified independently of $\gamma$, since the intratemporal labor-supply condition is free of wealth effects. The labor disutility scale, $\Psi$, has no independent economic content and is calibrated as a normalization so that average hours worked among employed prime-age men equal one third of the time endowment, consistent with a forty-hour workweek.

We calibrate the two parameters governing the health amenity, $\Xi$ and $\xi$, jointly. Following \citet{hall_jones_2007}, we choose the curvature parameter $\xi$ to match the income elasticity of health spending implied by cross-sectional and life-cycle variation in medical expenditures in ENASEM. We then discipline the scale parameter $\Xi$ by matching the value of a statistical life (VSL). We adopt a VSL of \$11.6 million in 2022 dollars for the United States.

Because $s(h)$ enters inside the GHH composite, the calibration of $\Xi$ interacts mechanically with the curvature parameter $\gamma$. We therefore solve for $(\Xi,\xi)$ taking $\gamma$ as given. Finally, we calibrate the discount factor, $\beta$, jointly with the rest of the life-cycle block to match an aggregate capital-output ratio of $3.0$.


\section{Data}\label{sec:data}

This section documents the empirical foundations of our quantitative
analysis. We rely on a combination of international databases and
national statistical sources to parameterize the model for the three
countries under study: Mexico, Costa Rica, and Panama. Together, these
countries offer a compelling comparative setting within Latin America:
they share institutional features---territorial proximity,
predominantly Spanish-speaking populations, and mixed public-private
health architectures---yet differ meaningfully in demographic
trajectories, fiscal capacity, health financing structures, and labour
market outcomes by gender. These contrasts make the three-country
framework particularly well suited for identifying how fiscal policy
interacts with health investment and gender dynamics within an
overlapping generations framework.

Our data compilation draws primarily from the WHO Global Health
Expenditure Database \citep{WHOGHED2024}, the OECD Health at a Glance
series \citep{OECDHaG2023}, the World Bank World Development
Indicators \citep{WorldBankWDI2024}, the Pan American Health
Organization platform \citep{PAHOHIA2024}, ECLAC and ILO joint labour
market reports \citep{ECLACILO2023}, and national tax policy summaries
\citep{PwCWWTS2024, OECDRevStats2024}. All monetary variables are
expressed in current USD unless otherwise indicated. Health expenditure
data refer to the 2022 reference year, the most recent for which
WHO-GHED provides complete cross-national data, while labour market
and fiscal indicators are drawn from 2022-2023 sources.

\subsection{Demographics and Health Status}\label{subsec:demog}

The three countries display heterogeneous demographic profiles that
shape life-cycle savings incentives and the public financing needs of
health and pension systems. Life expectancy at birth in 2023 was 75.1
years in Mexico, 81.0 years in Costa Rica, and 78.1 years in Panama
\citep{PAHOHIA2024, WHOGHO2024}. Costa Rica's longevity indicators are,
remarkably, on par with many high-income OECD economies, a reflection
of its longstanding single-payer insurance system and historically
strong primary care network. In all three countries, substantial gender
gaps in survival persist: females outlive males by approximately 6--7
years across the three settings, a pattern central to the gendered
saving and labour supply dynamics our model seeks to capture.

Fertility rates reveal a distinct divergence. Mexico's total fertility
rate was 1.82 births per woman in 2023, Costa Rica's stands at
1.51---below replacement level and among the lowest in Central
America---while Panama's is 2.30, reflecting relatively younger
demographic dynamics \citep{WorldBankWDI2024}. These trajectories imply
different paces of population ageing and, consequently, different
fiscal pressures on pension and health systems over the model horizon.

Non-communicable disease burdens are substantial in all three countries
and constitute a key driver of both public and out-of-pocket health
expenditures. Diabetes prevalence is particularly striking: Mexico
registers 16.9\% of its adult population living with diabetes---among
the highest rates globally---compared to 9.8\% in Costa Rica and
11.0\% in Panama \citep{OECDHaG2023}. These figures translate directly
into lifecycle health cost trajectories that are substantially higher
in Mexico than in the comparators, a pattern our model accommodates
through age- and health-state-dependent medical expenditure profiles.

Table~\ref{tab:demog} summarises the key demographic and health status
indicators that may serve as reference for parametrising the
demographic block of the model.

\begin{table}[htbp]
\centering
\caption{Demographic and Health Status Indicators}
\label{tab:demog}
\begin{tabular}{lccc}
\hline\hline
\textbf{Indicator} & \textbf{Mexico} & \textbf{Costa Rica} & \textbf{Panama} \\
\hline
Life expectancy at birth, total (2023)           & 75.1 & 81.0 & 78.1 \\
Life expectancy, male (2023)                     & 72.1 & 78.6 & 75.4 \\
Life expectancy, female (2023)                   & 78.3 & 83.5 & 80.9 \\
Fertility rate (births per woman, 2023)          & 1.82 & 1.51 & 2.30 \\
Infant mortality (per 1{,}000 live births, 2022) & 12.4 & 8.7  & 14.5 \\
Population aged 65+ (\% of total, 2023)          & 8.1  & 10.4 & 8.9  \\
Diabetes prevalence (\% of adults)               & 16.9 & 9.8  & 11.0 \\
\hline\hline
\end{tabular}
\begin{minipage}{\linewidth}
\smallskip
\footnotesize
\textit{Sources:} \citet{PAHOHIA2024}; \citet{WHOGHO2024};
\citet{WorldBankWDI2024}; \citet{OECDHaG2023}.
\end{minipage}
\end{table}

\subsection{Health Expenditure}\label{subsec:health_exp}

Health financing structures differ markedly across the three countries
and are critical for calibrating both the subsidy rate on medical
expenditures and the government budget constraint in the model. Current
health expenditure as a share of GDP in 2022 was approximately 5.5\%
in Mexico, 7.1\% in Costa Rica, and 7.8\% in Panama \citep{WHOGHED2024}.
These aggregate magnitudes, however, conceal substantial differences in
how expenditure is financed.

Costa Rica's health system is anchored by the \textit{Caja
Costarricense de Seguro Social} (CCSS), a compulsory social insurance
scheme that finances approximately 71\% of total health expenditure
through mandatory contributions; government and compulsory schemes
together account for close to 69\% of current health expenditure, while
out-of-pocket (OOP) spending represents 22.4\% of current health
expenditure (CHE)---equivalent to USD~220 per capita at current prices
\citep{OECDHaG2023, WHOGHO2024}. OOP spending as a share of final
household consumption was estimated at 2.6\% in 2022, reflecting the
strong financial protection that Costa Rica's system provides to
households.

Mexico presents a considerably different picture. Despite recent
reforms---including the phasing out of \textit{Seguro Popular} and the
introduction of IMSS-Bienestar---the share of health expenditure
financed out of pocket remained elevated, at approximately 41\% of
total health expenditure in 2021, down modestly from 44\% in 2010
\citep{WHOGHED2024}. This level substantially exceeds the WHO's
recommended threshold of 15--20\% for financial protection, implying
significant exposure to catastrophic health spending risk, particularly
among lower-income households. Mexico's per capita public health
spending---at roughly USD~626 PPP---is approximately half that of
Chile, Colombia, Costa Rica, or Panama \citep{OECDHaG2023}.

Finally, Panama's health financing is more mixed. Its public health system
encompasses the Ministry of Health (MINSA) and the Social Security Fund
(\textit{Caja de Seguro Social}, CSS), which together account for
approximately 72\% of current health expenditure. The OOP share hovers
around 26\% of CHE, placing Panama between Costa Rica and Mexico in
terms of household financial protection.

These financing structures may serve as empirical reference for
calibrating parameters related to the public financing of health care,
such as a potential subsidy rate on medical expenses borne by
households. Countries with
lower OOP ratios, Costa Rica most prominently, imply higher effective
subsidisation across the lifecycle, compressing the precautionary
savings motive relative to Mexico. The government budget constraint
must therefore accommodate not only demographic transitions but also
changes in the fiscal cost of these subsidy regimes.

Table~\ref{tab:health_exp} presents the health expenditure indicators
that may serve as reference for the calibration exercise.

\begin{table}[htbp]
\centering
\caption{Health Expenditure Indicators (Reference Year: 2022)}
\label{tab:health_exp}
\begin{tabular}{lccc}
\hline\hline
\textbf{Indicator} & \textbf{Mexico} & \textbf{Costa Rica} & \textbf{Panama} \\
\hline
Current health expenditure (\% of GDP)       & 5.5          & 7.1   & 7.8   \\
Public health expenditure (\% of CHE)        & $\approx$51  & $\approx$69 & $\approx$72 \\
Out-of-pocket expenditure (\% of CHE)        & $\approx$41$^{a}$ & 22.4 & $\approx$26 \\
OOP per capita, USD current                  & $\approx$253 & 220   & $\approx$290 \\
Total health exp.\ per capita, USD PPP       & $\approx$1{,}100 & 1{,}658 & $\approx$1{,}450 \\
OOP as \% of household final consumption     & $\approx$5.0 & 2.6   & $\approx$3.5 \\
\hline\hline
\end{tabular}
\begin{minipage}{\linewidth}
\smallskip
\footnotesize
\textit{Sources:} \citet{WHOGHED2024}; \citet{OECDHaG2023};
\citet{WHOGHO2024}. CHE = Current Health Expenditure; OOP =
Out-of-pocket; PPP = Purchasing Power Parity. Values marked
($\approx$) are approximations. $^{a}$~Mexico's OOP share
corresponds to the 2021 data vintage, the most recent available.
\end{minipage}
\end{table}

\subsection{Fiscal Parameters}\label{subsec:fiscal}

The fiscal block of the model is parameterised using statutory tax
rates and revenue aggregates from official sources. Three main fiscal
instruments enter the government's budget constraint: labour income
taxes ($\tau^{\ell}$), consumption taxes embedded in the value-added
tax (VAT), and social security contributions that finance both the
health and pension systems.

Consumption taxes vary considerably across the three countries. Mexico
imposes a standard VAT (\textit{Impuesto al Valor Agregado}, IVA) of
16\%, with exemptions covering basic food staples and medicines. Costa
Rica introduced its own IVA at a standard rate of 13\% following the
2018 fiscal reform (Law~9635), replacing an earlier general sales tax;
reduced rates of 4\% and 2\% apply to certain categories including
medicines and basic food. Panama's \textit{Impuesto de Transferencia de
Bienes Corporales Muebles y la Prestaci\'{o}n de Servicios} (ITBMS)
operates at a standard rate of 7\%, substantially lower than its
neighbours, with medicines and health services generally exempt
\citep{PwCWWTS2024}.

Progressive personal income tax schedules apply in all three countries,
with top marginal rates of 35\% in Mexico, 25\% in Costa Rica, and
25\% in Panama \citep{PwCWWTS2024}. These rates interact with social
security payroll contributions: total payroll contributions (employer
plus employee) amount to approximately 30\% of gross wages in Mexico
(IMSS), around 37\% in Costa Rica (CCSS), and roughly 19\% in Panama
(CSS) \citep{CCSS2023, PwCWWTS2024}. These contributions fund the
compulsory health insurance underpinning the social insurance schemes
described in Section~\ref{subsec:health_exp}.

Tax revenue as a share of GDP is estimated at approximately 17\% in
Mexico, 14\% in Costa Rica, and 12\% in Panama for 2022
\citep{OECDRevStats2024, WorldBankWDI2024}. Mexico's relatively higher
tax-to-GDP ratio conceals important structural weaknesses: a large
informal economy reduces the effective tax base, contributing to the
segmented access to social insurance that characterises its health
system. These parameters may serve as empirical reference for potential policy scenarios related to fiscal instruments such as labour income taxes, health subsidies, and pension transfers, particularly in the context of analysing government responses to demographic pressures.

Table~\ref{tab:fiscal} in Appendix~\ref{app:fiscal_table} presents the key fiscal indicators by country.

\subsection{Labour Market}\label{subsec:labour}

The labour market block of the model requires data on lifecycle
earnings profiles, participation rates by gender, the distribution of
income across skill groups, and the extent of informality. Gender
dynamics are of particular interest given the paper's focus: persistent
gender gaps in labour force participation and wages generate
differentiated savings trajectories and health expenditure patterns
that an OLG model with heterogeneous agents is well positioned to
capture.

Labour force participation rates reveal pronounced gender gaps across
all three countries. In 2022, male participation rates were 76.7\% in
Mexico, 72.4\% in Costa Rica, and 77.6\% in Panama, while female rates
stood at 45.4\%, 50.7\%, and 51.8\%, respectively \citep{ECLACILO2023,
ILOILOSTAT2023}. The male-to-female participation ratio was 1.69 in
Mexico, 1.43 in Costa Rica, and 1.50 in Panama. Mexico's gap is the
largest of the three and one of the widest in Latin America, a
structural feature linked to high informality, limited access to
childcare, and cultural norms around the sexual division of domestic
labour \citep{IMF2023, ECLAC2023}.

Gender wage gaps are significant in all three countries. Mexico's
gender wage gap was estimated at approximately 14\% of male earnings in
2021, slightly above the OECD average of 11.9\% \citep{IMF2023}. Costa
Rica and Panama face gaps in the range of 18--19\% of male wages
\citep{ILOILOSTAT2023, WEFGGR2024}. These gaps compound over the
lifecycle: women's lower lifetime earnings translate into lower pension
entitlements and a weaker capacity to self-insure against health
expenditure risk in old age, a mechanism central to the model's welfare
analysis.

Informality is a defining feature of all three labour markets.
Approximately 55\% of employed workers in Mexico operate informally,
compared to roughly 43\% in Costa Rica and 47\% in Panama
\citep{ECLACILO2023}. Workers in the informal sector typically lack
access to social-security-financed health insurance, implying that a
large fraction of the population must finance health care out of
pocket. 

Table~\ref{tab:labour} in Appendix~\ref{app:labour_table} summarises the main labour market indicators.

%% ====================================================================
%% SECTION 6 — POLICY COUNTERFACTUALS
%% ====================================================================
\section{Policy Counterfactuals}
\label{sec:policy}

\subsection{Framing the experiment}
\label{subsec:policy_framing}
 
We use the calibrated model to study how the fiscal pressures generated by demographic change can be partially absorbed by policy. Demographic change enters the model along two dimensions: a fall in population growth $n_p$ and a rise in age-conditional survival $\psi_{j+1}(\cdot)$ (longer life expectancy). Both raise the old-age dependency ratio $N^R/N^W$ and, given the pension closure rule~\eqref{eq:pension_contribution}, mechanically push the contribution rate $\tau^p$ upward. Both also reshape household savings, labor supply, and health investment, so the response of the equilibrium is not exhausted by the mechanical pension channel. The exercise has two steps. In the first one, we compute the new stationary equilibrium under the post-demographic-change values of $(n_p, \{\psi_{j+1}\})$, holding all policy instruments at their baseline values. This isolates the welfare and fiscal consequences of the demographic shift in the absence of any policy response. In the second one, we re-solve the model under alternative policy configurations and ask which combinations of instruments deliver smaller welfare losses, smaller debt accumulation, and smaller increases in $\tau^p$ relative to the baseline of inaction.
 
\subsection{Welfare criterion}
\label{subsec:policy_welfare}
The model has two skill types $\theta \in \{\theta_L, \theta_H\}$ and within-type heterogeneity in $(a, h, \eta)$. We report welfare at two levels of aggregation.
 
\paragraph{Newborn welfare.} For a newborn of type $\theta$, expected lifetime utility under policy regime $p$ is
\begin{equation}
\label{eq:welfare_newborn}
\mathcal{W}^p_1(\theta)
\;=\;
\mathbb{E}_1\!\left[\sum_{j=1}^{J} \beta^{j-1}
\left(\prod_{i=2}^{j} \psi_i(h_i^p)\right)
u\bigl(c_j^p,\, \ell_j^p,\, h_j^p\bigr)\right],
\end{equation}
where the expectation is over the path of $\eta$ realizations and the prefix superscript $p$ denotes policy-regime-conditional optimal paths. For each policy comparison we report the consumption-equivalent variation (CEV) of switching from regime $p$ to regime $p'$, defined by skill type as the constant proportional consumption increment $\Delta_\theta$ that solves
\begin{equation}
\label{eq:cev_newborn}
\mathbb{E}_1\!\left[\sum_{j=1}^{J} \beta^{j-1}
\left(\prod_{i=2}^{j} \psi_i(h_i^p)\right)
u\bigl((1 + \Delta_\theta)\, c_j^p,\, \ell_j^p,\, h_j^p\bigr)\right]
\;=\;
\mathcal{W}^{p'}_1(\theta).
\end{equation}


\paragraph{Aggregate welfare.} A scalar summary, weighted by the cohort weights:
\begin{equation}
\label{eq:welfare_agg}
\mathcal{W}^p \;=\; \sum_\theta \phi_1(\theta)\, \mathcal{W}^p_1(\theta),
\end{equation}
together with the population-weighted average CEV across types. We caution that the aggregate hides the bulk of the distributional content of the experiments, particularly the differential incidence by skill type that is the model's primary distributional contribution.
 
\subsection{Closure rule and what adjusts}
\label{subsec:policy_closure}
 
In all counterfactuals the pension contribution rate $\tau^p$ adjusts mechanically to balance the PAYG block, equation~\eqref{eq:pension_contribution}; the replacement rate $\kappa$, when not the policy lever under study, is held fixed. The general fiscal budget~\eqref{eq:gov_budget} closes with stationary government debt $B$ as the residual, equation~\eqref{eq:gov_stationary}. Tax rates $(\tau^c, \tau^\omega, \tau^k, \tau^m)$ and the public-spending share $g_y$ are exogenous. When one of these is the policy lever under study, we hold the others fixed and let $B$ absorb the residual change in revenue. We report alternative closures (one tax adjusts, $B$ held fixed) as robustness.

\subsection{Step 1: cost of inaction}
\label{subsec:policy_step1}

Under the baseline policy and post-shock demographics we report:
\begin{enumerate}
\item \textbf{Aggregate effects.} Equilibrium $(K, Y, C, M, w, r)$ relative to the pre-shock stationary equilibrium. Under longer life expectancy, lifecycle savings rise (precautionary motive against more retirement years), pushing $K$ up and $r$ down; under smaller $n_p$, the labor force shrinks relative to the capital stock, pushing $w$ up and $r$ down. Both forces compress $r$ and may push the economy toward dynamic inefficiency.
\item \textbf{Fiscal effects.} The dependency ratio $N^R/N^W$ and therefore $\tau^p$ rise mechanically. Capital and labor income tax bases shift; medical subsidy outlays $\tau^m \cdot M$ rise as healthier-but-older households invest more in $m$. We report stationary $B/Y$ before and after.
\item \textbf{Distributional effects.} CEV by type $\theta$ and by cohort age $j$ at the time of the shock. In the standard incidence pattern, the burden falls on currently-working middle-aged cohorts (longer expected retirement; higher $\tau^p$ for fewer of their working years) and is muted for retirees (already collecting) and far-future cohorts (born into the new steady state, who reoptimize fully).
\end{enumerate}
\subsection{Step 2: a menu of policy responses}
\label{subsec:policy_step2}
We organize the policy instruments into three groups.
 
\subsubsection{Pension-system levers: replacement rate vs. retirement age}
 
Both levers address the same fiscal pressure, but along different margins of the household problem. Cutting the replacement rate $\kappa$ is an \emph{intensive-margin} reform: retirees collect less per year. Raising $j_R$ is an \emph{extensive-margin} reform: retirees collect for fewer years. The two have very different incidence:
 
\begin{itemize}
\item \emph{Cutting $\kappa$} reduces transfers to retirees uniformly, raises private savings of pre-retirement cohorts (precautionary), and lowers $\tau^p$. Retirees with little wealth -- disproportionately of skill type $\theta_L$, who have lower lifetime productivity -- bear more of the cost. CEV is sharply negative for current retirees.
\item \emph{Raising $j_R$} extends the working life, which raises both $L$ and the pension tax base. It is welfare-improving in the model's aggregate utilitarian sense whenever the disutility of an extra working year is smaller than the welfare cost of the equivalent $\kappa$ cut. The incidence is again skewed: skill type $\theta_L$ faces a higher health-productivity penalty $\varrho(\theta_L)$, and so an additional working year has a higher utility cost for low-skill workers in poor health. The reform is therefore also distributionally regressive in a non-obvious way that the model captures.
\end{itemize}
 
\subsubsection{Tax-mix levers: $\tau^\omega$, $\tau^c$, $\tau^k$, $\tau^m$}
 
These instruments redistribute the fiscal burden across margins of household behavior:
 
\begin{itemize}
\item Raising $\tau^\omega$ distorts the labor-leisure FOC~\eqref{eq:cons_leisure}, reducing $L$. This is the textbook labor-tax distortion, magnified here because it stacks on top of the pension contribution $\tau^p$ that is already rising mechanically. Of the four taxes, $\tau^\omega$ is the most directly affected by the demographic shock and so has the smallest remaining margin to absorb fiscal pressure.
\item Raising $\tau^c$ is non-distortionary in the labor-leisure FOC (in the GHH structure here), but reduces the after-tax price of leisure relative to $\tau^\omega$, so the substitution effects differ. It is also broadly proportional in incidence across skill types.
\item Raising $\tau^k$ has the standard Diamond-style ambiguity: it distorts the savings decision via the Euler equation~\eqref{eq:euler}, but in a dynamically inefficient region (which demographic change pushes the economy toward), it can move the economy closer to the Golden Rule.
\item Adjusting $\tau^m$ -- the only lever with a non-zero subsidy interpretation -- works through four channels of the model and is the natural object to examine separately.
\end{itemize}
 
\subsection{Health subsidies and the four-channel mechanism}
\label{subsec:policy_health}

 
Adjusting $\tau^m$ is qualitatively different from the other tax instruments because medical spending $m$ has four distinct effects in the household problem, each operating through a different equation:
 
\begin{enumerate}
\item \emph{Productivity channel:} higher $h$ raises labor productivity $\nu_j(h, \eta; \theta) = e_j \exp(\theta + \eta - \varrho(\theta) h)$. This potentially can have heterogeneous effects across types.
\item \emph{Utility channel:} higher $h$ raises utility directly through the amenity term $s(h)$ in equation~\eqref{eq:utility}. This effect is symmetric across skill types.
\item \emph{Survival channel:} higher $h$ raises $\psi_{j+1}(h')$, extending expected life. This is the channel that interacts most subtly with the policy environment: health subsidies that raise life expectancy worsen the dependency ratio $N^R/N^W$, partially offsetting the fiscal benefit of any other reform. 
\item \emph{Resource cost:} medical spending consumes final goods at a rate $(1+\tau^m)m$ per unit of $h$ improvement (with $\tau^m < 0$ a subsidy), entering the resource constraint~\eqref{eq:goods_market} as part of $M$.
\end{enumerate}
 
The implication is that health subsidies are not a standalone fiscal instrument. They simultaneously alleviate fiscal pressure (productivity channel raising $L$ and tax bases) and exacerbate it (survival channel raising $N^R$). Whether the net fiscal impact is favorable is a quantitative question that the model is built to answer; we therefore study $\tau^m$ jointly with $\kappa$ or $j_R$ rather than in isolation.

\par In general, we use our model around the following three questions:
\begin{enumerate}
\item How much welfare does the economy lose from demographic change under inaction, and where does the loss fall (skill type, cohort age)?
\item Which single-instrument policy response delivers the smallest aggregate-welfare loss for a given target on $B/Y$ and $\tau^p$? Which is most regressive? Most progressive?
\item Are there policy combinations -- e.g., partial $\kappa$ cut combined with $\tau^m$ subsidy and a small rise in $j_R$ -- that dominate any single-instrument response?
\end{enumerate}
The model is built to answer all three. The third is the most demanding and the most policy-relevant: real demographic-change reforms in OECD countries combine all of these instruments, and the model's primary contribution is to quantify the trade-offs between them under endogenous mortality and skill heterogeneity for three Latin American countries.


%% ====================================================================
%% SECTION 7 — FORMALIZATION EXPERIMENT
%% ====================================================================
\section{Formalization Experiment}
\label{sec:formalization}

The policy counterfactuals in Section~\ref{sec:policy} study how fiscal instruments can absorb the pressures generated by demographic change within the formal-sector economy. This section asks a complementary question: what are the fiscal consequences of expanding the contributory base itself?

\subsection{Design}
\label{subsec:form_design}

Let the total working-age population be partitioned into formal workers $B$ and informal workers $C$, so that $A = B + C$. The baseline model is calibrated to population $B$: agents who contribute to the pension system, pay labor income tax $\tau^\omega$ and pension contributions $\tau^p$, and are covered by the contributory health system. Population $C$ operates outside this system---its members pay neither $\tau^\omega$ nor $\tau^p$ and are not covered by contributory health insurance.

The two populations differ systematically in skill composition. High-skilled workers ($\theta_H$) are overrepresented in $B$; low-skilled workers ($\theta_L$) are overrepresented in $C$. This sorting is not random---it reflects the equilibrium outcome of contribution costs, benefit valuations, enforcement, and institutional credibility documented in the formalization incentives literature \citep{WorldBankInformalizing2025, BoschMelguizoPages2013}.

The experiment proceeds as follows. We simulate a formalization program that transfers a fraction $\alpha \in (0, 1]$ of population $C$ into the contributory system. The newly formalized workers enter $B$ with their observed characteristics: income profiles, health trajectories, and survival probabilities drawn from national household survey data---ENOE for M\'exico, ENAHO for Costa Rica, and the Encuesta de Hogares for Panam\'a. The model is then re-solved for the new stationary equilibrium under the expanded population $B' = B + \alpha C$, holding all policy instruments at their baseline values.

\subsection{Competing Fiscal Effects}
\label{subsec:form_effects}

The formalization of informal workers generates two opposing fiscal forces.

\paragraph{Contribution-base effect.} The mass of contributing workers $N^W$ increases, while the mass of retirees $N^R$ rises only with a lag (newly formalized workers must accumulate contribution years before qualifying for pension benefits). In the short run, the dependency ratio $N^R/N^W$ falls, reducing the equilibrium pension contribution rate $\tau^p$ via equation~\eqref{eq:pension_contribution}. Total labor income tax revenue $\tau^\omega w L$ also rises as effective labor $L$ expands.

\paragraph{Expenditure-cost effect.} The newly formalized workers are predominantly low-skilled. They face earlier health depreciation (higher $\delta_j^h$ effective rates), lower labor productivity (lower $e_j$ profiles and higher health-productivity penalties $\varrho(\theta_L)$), and lower lifetime earnings. Their incorporation into the contributory health system raises per-capita public health expenditure. Over the lifecycle, their earlier health deterioration generates higher medical spending needs and, through the survival channel, potentially longer periods of pension receipt if health investments extend their lives.

\paragraph{Net fiscal impact.} The sign of the net effect depends on the relative magnitudes of these two forces, which vary by country. In M\'exico, where informality exceeds 55 percent and the informal population is heavily concentrated among low-skilled workers, the expenditure-cost effect is expected to be large. In Costa Rica, where informality is lower and the CCSS already provides broad coverage, the marginal fiscal cost of formalization may be smaller. The model quantifies these trade-offs by reporting, for each country and each formalization rate $\alpha$:

\begin{itemize}
    \item The change in the equilibrium pension contribution rate $\Delta \tau^p(\alpha)$.
    \item The change in stationary government debt $\Delta B/Y(\alpha)$.
    \item The consumption-equivalent variation by skill type, $\Delta_{\theta_L}(\alpha)$ and $\Delta_{\theta_H}(\alpha)$, for workers already in the formal sector---capturing the welfare externality that formalization imposes on incumbent formal workers through factor prices and pension contributions.
    \item The threshold formalization rate $\alpha^*$ at which the contribution-base effect exactly offsets the expenditure-cost effect on $\tau^p$.
\end{itemize}

\subsection{Interaction with Policy Counterfactuals}
\label{subsec:form_interaction}

The formalization experiment can be combined with the policy instruments studied in Section~\ref{sec:policy}. Two interactions are of particular interest.

First, a health subsidy ($\tau^m < 0$) targeted at newly formalized workers could partially offset their earlier health depreciation, raising their productivity and extending their productive working lives. This would strengthen the contribution-base effect relative to the expenditure-cost effect, potentially lowering $\alpha^*$. The model can quantify the fiscal cost of such a targeted subsidy against its indirect revenue gains through higher $L$ and lower $\tau^p$.

Second, a reduction in the pension vesting threshold $\bar{q}$ would allow newly formalized workers to qualify for pension benefits with fewer contribution years, raising their valuation of formal employment---the ``social security wealth'' channel emphasized by \citet{WorldBankInformalizing2025}. This is a demand-side complement to the supply-side formalization program: it makes formality more attractive to workers at the margin. The model traces the fiscal cost of earlier pension eligibility against the revenue gains from a voluntarily larger formal sector.

These joint experiments address the question of whether policy complementarities are quantitatively important in the context of formalization, and connect the model's fiscal analysis directly to the reform roadmap proposed by \citet{WorldBankInformalizing2025}.


%% ====================================================================
%% SECTION 8 — CONCLUSION
%% ====================================================================
\section{Conclusion}
\label{sec:conclusion}

This paper presents a fully specified analytical framework for studying the fiscal consequences of demographic aging in Latin America through the lens of health expenditure, endogenous longevity, and social security sustainability. The framework consists of a stochastic overlapping-generations model with continuous endogenous health capital, heterogeneous agents who differ by skill type, a production sector with endogenous factor prices, and a government that operates a pay-as-you-go pension system with an endogenous contribution rate. The model is calibrated for M\'exico, Costa Rica, and Panam\'a using data from the Global Burden of Disease, the OECD Taxing Wages series, the WHO Global Health Expenditure Database, and national household surveys.

The model's central analytical contribution is the identification of a fiscal amplification loop that links health investment to pension sustainability through the longevity channel. Health investment raises the health stock, which improves survival probabilities and extends working lives---strengthening the life-cycle savings motive, raising aggregate capital, and compressing the equilibrium interest rate. But longer lives also raise the dependency ratio $N^R/N^W$, mechanically increasing the pension contribution rate $\tau^p = \kappa \cdot N^R/N^W$ and widening the tax wedge on formal labor income. The net fiscal effect of improved population health is therefore ambiguous in sign. This ambiguity cannot be resolved by partial-equilibrium analysis: it requires the calibrated general-equilibrium model that this paper builds, in which health-driven changes in aggregate savings move factor prices and feed back into every household's decision problem.

The paper develops two complementary experimental frameworks. The policy counterfactual battery (Section~\ref{sec:policy}) studies how fiscal instruments---pension parameters ($\kappa$, $j_R$), tax rates ($\tau^\omega$, $\tau^c$, $\tau^k$), and health subsidies ($\tau^m$)---can absorb the fiscal pressures generated by demographic change. The four-channel mechanism through which health subsidies operate---productivity, utility, survival, and resource cost---makes $\tau^m$ qualitatively different from the other fiscal instruments and motivates its joint analysis with pension reforms. The formalization experiment (Section~\ref{sec:formalization}) studies the fiscal consequences of expanding the contributory base by incorporating informal workers with their observed income and health profiles, identifying the threshold formalization rate $\alpha^*$ at which the contribution-base effect dominates the expenditure-cost effect. Both frameworks are designed to evaluate policy complementarities---whether combinations of instruments dominate any single-instrument response---which is the most policy-relevant and computationally demanding question the model addresses.

The model's architecture supports a natural extension toward gender-differentiated fiscal analysis. By expanding the agent-type space from two skill types to four---$(\theta_L^M, \theta_H^M, \theta_L^F, \theta_H^F)$---calibrated with sex-specific earnings trajectories and the survival differentials documented in the data section, the model can decompose fiscal sustainability pressures by sex within the existing equilibrium framework. This calibration-based approach lets observed income profiles and survival probabilities carry the information that a structural household bargaining model would generate endogenously, at a fraction of the computational cost.

The remaining phases of the project are organized as follows. By the IDB Second Seminar (week of June 8, 2026), baseline results and preliminary calibrations will be reported for the three countries, and initial policy counterfactuals---including the cost-of-inaction benchmark and at least the pension-lever and health-subsidy experiments---will be presented. The formalization experiment with observed income profiles and the gender-differentiated decomposition will follow in the July--August period. The final synthesis, including all counterfactual results, the complete policy analysis, replication code, and datasets, is scheduled for submission by September 4, 2026.

The principle that organizes the analysis is straightforward: the question facing Latin American health systems is not whether to spend more, but \textit{when} to spend, through \textit{which institutional channels}, and with what consequences for the contributory architecture on which fiscal sustainability depends.


%% ====================================================================
%% APPENDICES
%% ====================================================================
\begin{appendices}

\section{Fiscal Gap Methodology}
\label{app:fiscal_gap}

The fiscal gap for country $c$ in year $t$ is defined as:
\begin{equation}
\label{eq:fiscal_gap}
G_{c,t} = H_{c,t}^{\text{est}} - H_{c,t}^{\text{actual}}
\end{equation}
where $H_{c,t}^{\text{est}}$ denotes estimated health expenditure requirements and $H_{c,t}^{\text{actual}}$ denotes actual public health expenditure, both expressed as shares of GDP. Estimated requirements are constructed by applying a base-year age-expenditure profile $h(j)$ to the projected population structure:
\begin{equation}
\label{eq:fiscal_gap_est}
H_{c,t}^{\text{est}} = \frac{\sum_j h(j) \cdot N_{j,c,t}}{Y_{c,t}}
\end{equation}
where $N_{j,c,t}$ is the population of age group $j$ in country $c$ at time $t$, and $Y_{c,t}$ is nominal GDP. The base-year expenditure profile $h(j)$ is estimated from national health accounts microdata for the reference year.


\section{Supplementary Tables and Data}
\label{app:supplementary}

\subsection{Fiscal Indicators}
\label{app:fiscal_table}

\begin{table}[H]
\centering
\caption{Fiscal Indicators}
\label{tab:fiscal}
\begin{tabular}{lccc}
\hline\hline
\textbf{Indicator} & \textbf{Mexico} & \textbf{Costa Rica} & \textbf{Panama} \\
\hline
Standard VAT rate                                  & 16\% (IVA) & 13\% (IVA) & 7\% (ITBMS) \\
Top personal income tax rate (2023)                & 35\%  & 25\%  & 25\%  \\
Social security contributions (employer + employee) & $\approx$30\% & $\approx$37\% & $\approx$19\% \\
General government revenue (\% of GDP, 2022)       & $\approx$23 & $\approx$24 & $\approx$21 \\
Tax revenue (\% of GDP, 2022)                      & $\approx$17 & $\approx$14 & $\approx$12 \\
\hline\hline
\end{tabular}
\begin{minipage}{\linewidth}
\smallskip
\footnotesize
\textit{Sources:} \citet{PwCWWTS2024}; \citet{OECDRevStats2024};
\citet{WorldBankWDI2024}; \citet{CCSS2023}.
\end{minipage}
\end{table}

\subsection{Labour Market Indicators}
\label{app:labour_table}

\begin{table}[H]
\centering
\caption{Labour Market Indicators (Reference Year: 2022)}
\label{tab:labour}
\begin{tabular}{lccc}
\hline\hline
\textbf{Indicator} & \textbf{Mexico} & \textbf{Costa Rica} & \textbf{Panama} \\
\hline
Male labour force participation rate (\%)    & 76.7        & 72.4        & 77.6        \\
Female labour force participation rate (\%)  & 45.4        & 50.7        & 51.8        \\
Gender participation gap (male/female ratio) & 1.69        & 1.43        & 1.50        \\
Gender wage gap (\% of male wage)            & $\approx$14 & $\approx$18 & $\approx$19 \\
Informality rate (\% of employed)            & $\approx$55 & $\approx$43 & $\approx$47 \\
Unemployment rate (\%)                       & 3.3         & 11.2        & 7.7         \\
Gini coefficient                             & 0.43        & 0.48        & 0.49        \\
\hline\hline
\end{tabular}
\begin{minipage}{\linewidth}
\smallskip
\footnotesize
\textit{Sources:} \citet{ILOILOSTAT2023}; \citet{ECLACILO2023};
\citet{IMF2023}; \citet{WEFGGR2024}; \citet{WorldBankWDI2024};
\citet{WorldBankPovCalNet2022}.
\end{minipage}
\end{table}

\subsection{Life Table for Mexico}
\label{app:survival}

The following life table corresponds to Mexico and is used to
calibrate the age-specific survival probabilities $\psi_{j}$ in the
model. Data are drawn from \citet{UISP2019}.

\begin{table}[H]
\centering
\caption{Life Table for Mexico: Mortality and Survival Probabilities}
\label{tab:survival_mx}
\begin{tabular}{lcc}
\hline\hline
\textbf{Age group} & \textbf{Probability of death} & \textbf{Survival probability} \\
\hline
$<$ 1 year   & 0.011576 & 0.9884 \\
1--4         & 0.002282 & 0.9977 \\
5--9         & 0.001132 & 0.9989 \\
10--14       & 0.001470 & 0.9985 \\
15--19       & 0.003858 & 0.9961 \\
20--24       & 0.005897 & 0.9941 \\
25--29       & 0.006550 & 0.9935 \\
30--34       & 0.007966 & 0.9920 \\
35--39       & 0.011217 & 0.9888 \\
40--44       & 0.013710 & 0.9863 \\
45--49       & 0.019216 & 0.9808 \\
50--54       & 0.028413 & 0.9716 \\
55--59       & 0.042568 & 0.9574 \\
60--64       & 0.063162 & 0.9368 \\
65--69       & 0.093160 & 0.9068 \\
70--74       & 0.137294 & 0.8627 \\
75--79       & 0.208926 & 0.7911 \\
80--84       & 0.312093 & 0.6879 \\
85+          & 1.000000 & 0.0000 \\
\hline\hline
\end{tabular}
\begin{minipage}{\linewidth}
\smallskip
\footnotesize
\textit{Source:} \citet{UISP2019}. Net probability of death and
survival probability by age group for the total population of Mexico.
\end{minipage}
\end{table}

\end{appendices}


%% ====================================================================
%% BIBLIOGRAPHY
%% ====================================================================
\bibliography{references}

\end{document}